\SourcePage{168}{171}
\section{Time-independent perturbations}

An exact solution of the Schrödinger equation can be found in only a
comparatively small number of the simplest cases. Most problems in quantum
mechanics lead to equations too complicated for exact solution. Frequently,
however, a problem contains quantities of different orders of magnitude,
including small quantities whose omission simplifies the problem enough to
make an exact solution possible. The first step is then to solve this
simplified problem exactly; the second is to calculate approximately the
corrections due to the small terms that were omitted. The general method for
calculating these corrections is called \emph{perturbation theory}.

Suppose that the Hamiltonian of the physical system is
\[
  \hat H=\hat H_0+\hat V,
\]
where $\hat V$ is a small correction, the \emph{perturbation}, to the
``unperturbed'' operator $\hat H_0$. In \secref{38} and \secref{39} we consider
perturbations $\hat V$ with no explicit time dependence; the same is assumed
for $\hat H_0$. The conditions under which $\hat V$ may be regarded as small
in comparison with $\hat H_0$ will be established below.

For the discrete spectrum, the problem may be stated as follows. The
eigenfunctions $\psi^{(0)}$ and eigenvalues of $\hat H_0$ are assumed known;
that is, the exact solutions of
\begin{equation}
  \hat H_0\psi^{(0)}=E^{(0)}\psi^{(0)}
  \tag{38.1}
\end{equation}
are known. Approximate solutions are required for
\begin{equation}
  \hat H\psi=(\hat H_0+\hat V)\psi=E\psi,
  \tag{38.2}
\end{equation}
\SourcePage{169}{172}
namely approximate expressions for the eigenfunctions $\psi_n$ and
eigenvalues $E_n$ of the perturbed operator $\hat H$.

Throughout this section all eigenvalues of $\hat H_0$ will be assumed
nondegenerate. To simplify the derivation, we first suppose that the energy
levels form a discrete spectrum only.

It is convenient to use matrix form from the outset. Expand the required
function $\psi$ in the functions $\psi_m^{(0)}$:
\begin{equation}
  \psi=\sum_m c_m\psi_m^{(0)}.
  \tag{38.3}
\end{equation}
Substitution in (38.2) gives
\[
  \sum_m c_m(E_m^{(0)}+\hat V)\psi_m^{(0)}
  =\sum_m c_mE\psi_m^{(0)}.
\]
Multiplying by $\psi_k^{(0)*}$ and integrating, we find
\begin{equation}
  (E-E_k^{(0)})c_k=\sum_mV_{km}c_m.
  \tag{38.4}
\end{equation}
Here $V_{km}$ is the matrix of $\hat V$ in the unperturbed functions:
\begin{equation}
  V_{km}=\int\psi_k^{(0)*}\hat V\psi_m^{(0)}\,dq.
  \tag{38.5}
\end{equation}
Seek the coefficients and energy as series
\[
  E=E^{(0)}+E^{(1)}+E^{(2)}+\ldots,\qquad
  c_m=c_m^{(0)}+c_m^{(1)}+c_m^{(2)}+\ldots,
\]
where $E^{(1)},c_m^{(1)}$ have the same order of smallness as $\hat V$,
$E^{(2)},c_m^{(2)}$ are of second order, and so forth.

To calculate the corrections to the $n$th eigenvalue and eigenfunction, set
$c_n^{(0)}=1$ and $c_m^{(0)}=0$ for $m\ne n$. For the first approximation,
put $E=E_n^{(0)}+E_n^{(1)}$ and
$c_k=c_k^{(0)}+c_k^{(1)}$ in (38.4), retaining first-order terms only. The
equation with $k=n$ gives
\begin{equation}
  E_n^{(1)}=V_{nn}=\int\psi_n^{(0)*}\hat V\psi_n^{(0)}\,dq.
  \tag{38.6}
\end{equation}
\SourcePage{170}{173}
Thus the first-order correction to $E_n^{(0)}$ is the mean value of the
perturbation in the state $\psi_n^{(0)}$.

For $k\ne n$, equation (38.4) gives
\begin{equation}
  c_k^{(1)}=\frac{V_{kn}}{E_n^{(0)}-E_k^{(0)}},\qquad k\ne n.
  \tag{38.7}
\end{equation}
The coefficient $c_n^{(1)}$ remains arbitrary and must be chosen so that
$\psi_n=\psi_n^{(0)}+\psi_n^{(1)}$ is normalized through first order. This
requires $c_n^{(1)}=0$. Indeed,
\begin{equation}
  \psi_n^{(1)}=\sum_m'\frac{V_{mn}}{E_n^{(0)}-E_m^{(0)}}\psi_m^{(0)}
  \tag{38.8}
\end{equation}
(the prime means that $m=n$ is omitted) is orthogonal to $\psi_n^{(0)}$;
hence the integral of $|\psi_n^{(0)}+\psi_n^{(1)}|^2$ differs from unity only
by a second-order quantity.

Formula (38.8) gives the first-order correction to the wave functions. It also
shows the condition for applicability of the method:
\begin{equation}
  |V_{mn}|\ll|E_n^{(0)}-E_m^{(0)}|.
  \tag{38.9}
\end{equation}
Thus perturbation matrix elements must be small compared with the
corresponding separations between unperturbed energy levels.

To find the second-order correction to $E_n^{(0)}$, substitute
$E=E_n^{(0)}+E_n^{(1)}+E_n^{(2)}$ and
$c_k=c_k^{(0)}+c_k^{(1)}+c_k^{(2)}$ in (38.4), and retain second-order terms.
The equation with $k=n$ gives
\[
  E_n^{(2)}c_n^0=\sum_m'V_{nm}c_m^{(1)},
\]
and consequently
\begin{equation}
  E_n^{(2)}=\sum_m'\frac{|V_{mn}|^2}{E_n^{(0)}-E_m^{(0)}}.
  \tag{38.10}
\end{equation}
Here (38.7) was used, together with Hermiticity,
$V_{mn}=V_{nm}^*$.
\SourcePage{171}{174}
The second-order correction to the ground-state energy is always negative:
if $E_n^{(0)}$ is the lowest eigenvalue, every term in (38.10) is negative.
Higher approximations may be calculated in the same manner.

These results extend directly to an operator $\hat H_0$ that also has a
continuous spectrum, while the perturbed state itself still belongs to the
discrete spectrum. The appropriate continuous-spectrum integrals are simply
added to the discrete sums. Denote continuous-spectrum states by an index
$\nu$ ranging continuously; $\nu$ stands for the complete set of quantities
needed to specify a state. If these states are degenerate, as they nearly
always are, energy alone does not specify one.\footnote{The wave functions
$\psi_\nu^{(0)}$ must then be normalized to delta functions of the quantities
$\nu$.} For example, (38.8) becomes
\begin{equation}
  \psi_n^{(1)}=\sum_m'\frac{V_{mn}}{E_n^{(0)}-E_m^{(0)}}\psi_m^{(0)}
  +\int\frac{V_{\nu n}}{E_n^{(0)}-E_\nu^{(0)}}\psi_\nu^{(0)}\,d\nu.
  \tag{38.11}
\end{equation}

It is also useful to give the perturbed matrix elements of a physical
quantity $f$, calculated through first order with
$\psi_n=\psi_n^{(0)}+\psi_n^{(1)}$ and (38.8):
\begin{equation}
  f_{nm}=f_{nm}^{(0)}
  +\sum_k'\frac{V_{nk}f_{km}^{(0)}}{E_n^{(0)}-E_k^{(0)}}
  +\sum_k'\frac{V_{km}f_{nk}^{(0)}}{E_m^{(0)}-E_k^{(0)}}.
  \tag{38.12}
\end{equation}
In the first sum $k\ne n$, and in the second $k\ne m$.

\begin{center}\ProblemHeading\end{center}

\textbf{1.} Find the second-order correction to the eigenfunctions.

\SolutionHeading The coefficients $c_k^{(2)}$ for $k\ne n$ are obtained from
(38.4) with $k\ne n$, written through second order; $c_n^{(2)}$ is chosen so
that $\psi_n=\psi_n^{(0)}+\psi_n^{(1)}+\psi_n^{(2)}$ is
\SourcePage{172}{175}
normalized through second order. The result is
\[
 \begin{aligned}
  \psi_n^{(2)}={}&
  \sum_m'\sum_k'\frac{V_{mk}V_{kn}}{\hbar^2\omega_{nk}\omega_{nm}}
      \psi_m^{(0)}
  -\sum_m'\frac{V_{nn}V_{mn}}{\hbar^2\omega_{nm}^2}\psi_m^{(0)}\\
  &-\frac{\psi_n^{(0)}}2
  \sum_m'\frac{|V_{mn}|^2}{\hbar^2\omega_{nm}^2},
 \end{aligned}
\]
where
\[
  \omega_{nm}=\frac1\hbar(E_n^{(0)}-E_m^{(0)}).
\]

\textbf{2.} Find the third-order correction to the energy eigenvalues.

\SolutionHeading Retaining the third-order terms in (38.4) with $k=n$ gives
\[
  E_n^{(3)}=\sum_k'\sum_m'
  \frac{V_{nm}V_{mk}V_{kn}}{\hbar^2\omega_{mn}\omega_{kn}}
  -V_{nn}\sum_m'\frac{|V_{nm}|^2}{\hbar^2\omega_{mn}^2}.
\]

\textbf{3.} Find the energy levels of the anharmonic linear oscillator with
Hamiltonian
\[
  \hat H=\frac{\hat p^2}{2m}+\frac{m\omega^2x^2}{2}+\alpha x^3+\beta x^4.
\]

\SolutionHeading Matrix elements of $x^3$ and $x^4$ follow directly from the
rule for matrix multiplication and the elements of $x$ in (23.4). The nonzero
elements of $x^3$ are
\[
  (x^3)_{n-3,n}=(x^3)_{n,n-3}
  =\left(\frac\hbar{m\omega}\right)^{3/2}
   \sqrt{\frac{n(n-1)(n-2)}8},
\]
\[
  (x^3)_{n-1,n}=(x^3)_{n,n-1}
  =\left(\frac\hbar{m\omega}\right)^{3/2}
   \sqrt{\frac{9n^3}{8}}.
\]
There are no diagonal elements in this matrix, so the term $\alpha x^3$
gives no first-order correction when regarded as a perturbation of the
harmonic oscillator. Its second-order correction has the same order as the
first-order correction from $\beta x^4$. The diagonal elements of $x^4$ are
\[
  (x^4)_{n,n}=\left(\frac\hbar{m\omega}\right)^2
  \frac34(2n^2+2n+1).
\]
Equations (38.6) and (38.10) then give
\[
 \begin{aligned}
  E_n={}&\hbar\omega\left(n+\frac12\right)
  -\frac{15}{4}\frac{\alpha^2}{\hbar\omega}
   \left(\frac\hbar{m\omega}\right)^3
   \left(n^2+n+\frac{11}{30}\right)\\
  &+\frac32\beta\left(\frac\hbar{m\omega}\right)^2
   \left(n^2+n+\frac12\right).
 \end{aligned}
\]

\textbf{4.} An infinite spherical potential well undergoes a small
volume-preserving deformation into a slightly prolate or oblate ellipsoid of
revolution with semiaxes $a=b$ and $c$. Find the splitting of its particle
energy levels (\emph{A.~B. Migdal}, 1959).
\SourcePage{173}{176}
\SolutionHeading The boundary equation
\[
  \frac{x^2+y^2}{a^2}+\frac{z^2}{c^2}=1
\]
becomes the sphere $x^2+y^2+z^2=R^2$ under
$x\to ax/R$, $y\to ay/R$, $z\to cz/R$. The particle Hamiltonian
$\hat H=\hat{\mathbf p}^2/2M=-\hbar^2\Delta/2M$ ($M$ is the mass; energy is
measured from the bottom of the well) becomes $\hat H=\hat H_0+\hat V$, with
\[
  \hat H_0=-\frac{\hbar^2}{2M}\Delta,\qquad
  \hat V=-\frac{\hbar^2}{2M}\left[
   \left(\frac{R^2}{a^2}-1\right)
      \left(\frac{\partial^2}{\partial x^2}+\frac{\partial^2}{\partial y^2}\right)
   +\left(\frac{R^2}{c^2}-1\right)\frac{\partial^2}{\partial z^2}
  \right].
\]
The ellipsoidal-well problem is thereby reduced to the spherical-well
problem. If the ellipsoid differs only slightly from a sphere of radius
$R=(a^2c)^{1/3}$, $\hat V$ is a small perturbation. Introducing the
``ellipticity'' $\beta$, $|\beta|\ll1$, by
\[
  a\approx R\left(1-\frac\beta3\right),\qquad
  c\approx R\left(1+\frac{2\beta}3\right),
\]
we obtain
\[
  \hat V=\frac\beta{3M}(\hat{\mathbf p}^2-3\hat p_z^2).
\]
In first order, the change from the spherical-well levels is
\[
  \Delta E_{nlm}=E_{nlm}-E_{nl}^{(0)}=\langle nlm|V|nlm\rangle.
\]
Here $l,m$ are the angular momentum and its projection on the ellipsoid axis,
and $n$ numbers the spherical-well levels for fixed $l$, which do not depend
on $m$. The expression $\mathbf p^2-3p_z^2$ is the $zz$ component of the
traceless irreducible tensor $\delta_{ik}\mathbf p^2-3p_ip_k$. By (107.2)
and (107.6), $\langle nlm|V|nlm\rangle$ is proportional to
$(-1)^m\begin{pmatrix}l&2&l\\-m&0&m\end{pmatrix}$, and hence
\[
  \langle nlm|V|nlm\rangle
  =\left(1-\frac{3m^2}{l(l+1)}\right)\langle nl0|V|nl0\rangle.
\]
(The table of $3j$ symbols is on p.~530.) We next write
\[
 \begin{aligned}
  \langle nl0|V|nl0\rangle
  &=\frac23\beta E_{nl}^{(0)}
    +\beta\frac{\hbar^2}{M}
      \left\langle nl0\left|\frac{\partial^2}{\partial z^2}\right|nl0\right\rangle\\
  &=\frac23\beta E_{nl}^{(0)}
    -\frac{\beta\hbar^2}{M}\int
      \left|\frac{\partial\psi_{nl0}}{\partial z}\right|^2r^2\,dr\,do.
 \end{aligned}
\]
The Schrödinger equation
$\hat H_0\psi_{nlm}=E_{nl}^{(0)}\psi_{nlm}$ was used in the first term, and
the second was integrated by parts.
\SourcePage{174}{177}
Using (28.11) for $Y_{l0}$, the derivative of
$\psi_{nl0}=R_{nl}(r)Y_{l0}(\theta,\varphi)$ is
\[
 \begin{aligned}
  \frac\partial{\partial z}\psi_{nl0}
  &=\left(\cos\theta\frac\partial{\partial r}
    -\frac{\sin\theta}{r}\frac\partial{\partial\theta}\right)\psi_{nl0}\\
  &=-\frac{i(l+1)}{[4(l+1)^2-1]^{1/2}}
    \left(R_{nl}'-\frac lrR_{nl}\right)Y_{l+1,0}\\
  &\quad+\frac{il}{[4l^2-1]^{1/2}}
    \left(R_{nl}'+\frac{l+1}{r}R_{nl}\right)Y_{l-1,0}.
 \end{aligned}
\]
The radial integrals are evaluated from
\[
  \int_0^\infty R_{nl}R_{nl}'r\,dr
  =-\frac12\int_0^\infty R_{nl}^2\,dr,
\]
\[
  \int_0^\infty R_{nl}'^{,2}r^2\,dr
  =\frac{2M}{\hbar^2}E_{nl}^{(0)}
   -l(l+1)\int_0^\infty R_{nl}^2\,dr,
\]
by integration by parts and the radial Schrödinger equation (33.3),
\[
  R_{nl}''+\frac2rR_{nl}'-\frac{l(l+1)}{r^2}R_{nl}
  =-\frac{2M}{\hbar^2}E_{nl}^{(0)}.
\]
Terms containing integrals of $R_{nl}^2$ cancel, leaving
\[
  \Delta E_{nlm}=4\beta\frac{l(l+1)}{(2l-1)(2l+3)}
  \left[\frac{m^2}{l(l+1)}-\frac13\right]E_{nl}^{(0)}.
\]
Finally,
\[
  \frac1{2l+1}\sum_{m=-l}^lE_{nlm}=E_{nl}^{(0)},
\]
so the centre of gravity of the multiplet is unchanged.
